\newpage
\chapter{Performance}
The purpose of this function is to calculate the performances for a certain condition and handle the needed calculations for the following functions.
The calculation is based on the following assumptions: 
\begin{itemize}
	\item the combustion is in chemical equilibrium
	\item infinite area combustor
	\item frozen equilibrium at throat section
\end{itemize}
To perform this calculations NASA CEA \cite{NASACEA} and its Python wrapper
RocketCEA \cite{RocketCEA} were used.
\section{Functions}
The following functions are available in \texit{Performance//performance\_singlepoint.py}.
\begin{verbatim}
calculate_performance(Ainj, Aport, Ab, eps, ptank, Ttank, pc, CD,
                          a, n, rho_fuel, oxidizer, fuel, pamb=0.0, gamma0=1.3)
-> return (p_inj, mdot_ox, mdot_fuel, mdot, Gox, r, MR, Tc, MW, gamma, eps_out, cs,
            CF_vac, CF, Ivac, Is, flag_performance)
\end{verbatim}
\hspace{2em} Calculates the line losses and the injected mass flow, uses it to get oxidizer mass flux for the regression rate
\begin{equation}
	r=aG_{ox}^n
\end{equation}
to calculate the fuel mass flow
\begin{equation}
	\dot{m}_{fuel}=\rho_f A_b r
\end{equation}
from it the mixture ratio can be calculated as
\begin{equation}
	MR=\frac{\dot{m}_{ox}}{\dot{m}_{f}}
\end{equation}
With that we have the three needed inputs for the CEA calculations (chamber pressure, mixture ratio and expansion ratio).\\
It is possible to ask for the adapted expansion ratio passing \verb|eps="adapt"|, and it will be calculated as
\begin{equation}
	\varepsilon = \frac{\Gamma}{f(\frac{p_{amb}}{p_c})}
\end{equation}
with
\begin{align*}
	\Gamma&=\sqrt{\gamma (\frac{2}{\gamma +1})^\frac{\gamma +1}{\gamma -1}}\\
	f(\frac{p_{amb}}{p_1})&=\sqrt{\frac{2\gamma}{\gamma -1}((\frac{p_{amb}}{p_1})^\frac{2}{\gamma}-(\frac{p_{amb}}{p_1})^\frac{\gamma +1}{\gamma})}
\end{align*}

\begin{verbatim}
pressure_fun(Ainj, Aport, At, Ab, eps, ptank, Ttank, pc, CD, a, n, rho_fuel, 
			oxidizer, fuel, pamb=0.0, gamma0=1.3)
-> return Fpc
\end{verbatim}
\hspace{2em} From the definition
\begin{equation}
	c^*=\frac{p_cA_t}{\dot{m}}
\end{equation}
and we use it to calculate the pressure function to find the steady state pressure
for the optimisation:
\begin{equation}
	F(p_c)=\frac{\dot{m}c^*}{A_t} - pc
	\label{eq:PressFunction}
\end{equation}







